The quotient Cartier criterion also has the following version with a
finite-image endomorphism. Let $\alpha_1,\ldots,\alpha_r$ be commuting
endomorphisms of $J$ defined over $\mathbb Q$, commuting with
$\delta=\langle2\rangle$, and suppose that every
$\alpha_iJ(\mathbb Q)$ is finite. Replace the space in
Lemma~\ref{lem:quotient-Cartier} by
\[
 U=(\mathcal C-1)\sum_i\alpha_i^*H^0(Y_{\mathbb F_2},\Omega^1).
\]
The same conclusion holds if its restrictions have rank $d$.
Indeed, $\alpha_i[D-C]$ is torsion: write it as
$\alpha_i[D-dc_0]-\alpha_i[C-dc_0]$ and use respectively the finite-image
hypothesis and Manin--Drinfeld. Lemma~\ref{lem:dyadic-Hecke} annihilates
these reduction-kernel classes after applying $T_2-\delta$.
Use the product Abel map with components
$(T_2-\delta)\alpha_i[E-dc_0]$ in the formal-immersion argument.

For the curve of Lemma~\ref{lem:C65-genus41-model}, take
$\alpha_T=\tau_T-1$ and $\alpha_U=\tau_U-1$.
The only positive-rank newform factor at level dividing $65$ has trivial
diamond character. Both endomorphisms therefore have finite image on
$J(\mathbb Q)$. The two special-fibre lifts may each be changed by
$w\mapsto w+1$; both choices are diamond operators, so the finite-image
assertion holds for the explicit lifts used in the computation.

The same proof can be applied to a specified set of cuspidal reduction
divisors: surjectivity at each member of that set excludes noncuspidal
rational divisors whose reductions belong to it. This permits the
coefficient restrictions imposed by a modular quotient.
